\documentclass[12pt]{article}
\usepackage{ficha}

\begin{document}

\ficha{Sólidos Geométricos e Planificações}

\instrucoes{Resolve os exercícios. As figuras são esquemáticas e não estão à escala.}

\secao{Poliedros e Não Poliedros}{Classifica cada sólido como poliedro ou não poliedro e, no caso de poliedro, indica o seu nome.}
\begin{exercicios}
  \item Uma lata de refrigerante.
  \item Um dado.
  \item Uma bola de futebol.
  \item Uma caixa de sapatos.
  \item Um chapéu de aniversário (cone).
  \item Uma pirâmide do Egito.
  \item Uma tenda com duas faces triangulares e três retangulares.
  \item Um lápis afiado de secção hexagonal (considera o corpo do lápis).
  \item O que distingue um poliedro de um não poliedro?
  \item O que distingue um prisma de uma pirâmide?
\end{exercicios}

\secao{Faces, Vértices e Arestas}{Completa a tabela e verifica a relação de Euler $F + V = A + 2$.}

\begin{center}
\begin{tabular}{|l|c|c|c|c|}\hline
  Poliedro & Faces ($F$) & Vértices ($V$) & Arestas ($A$) & $F + V - A$ \\\hline
  Cubo & & & & \\\hline
  Prisma triangular & & & & \\\hline
  Prisma pentagonal & & & & \\\hline
  Pirâmide quadrangular & & & & \\\hline
  Pirâmide hexagonal & & & & \\\hline
  Octaedro & & & & \\\hline
\end{tabular}
\end{center}

\begin{exercicios}
  \item Quantas faces, vértices e arestas tem um prisma cuja base tem $n$ lados?
  \item Quantas faces, vértices e arestas tem uma pirâmide cuja base tem $n$ lados?
  \item Um prisma tem 18 arestas. Qual é o polígono da base?
  \item Uma pirâmide tem 12 arestas. Qual é o polígono da base?
  \item Um poliedro tem 8 faces e 12 vértices. Quantas arestas tem?
  \item Existe algum prisma com 15 arestas? Justifica.
\end{exercicios}

\secao{Planificações}{Observa as planificações e identifica o sólido correspondente.}

\begin{center}
\begin{tikzpicture}[scale=0.55]
  % (a) cubo
  \foreach \x in {0,1,2,3} \draw (\x,1) rectangle (\x+1,2);
  \draw (1,0) rectangle (2,1); \draw (1,2) rectangle (2,3);
  \node at (2,-0.8) {(a)};
  % (b) prisma triangular
  \begin{scope}[xshift=6cm]
    \foreach \x in {0,1.5,3} \draw (\x,0.5) rectangle (\x+1.5,2.5);
    \draw (1.5,2.5) -- (3,2.5) -- (2.25,3.8) -- cycle;
    \draw (1.5,0.5) -- (3,0.5) -- (2.25,-0.8) -- cycle;
    \node at (2.25,-1.6) {(b)};
  \end{scope}
  % (c) pirâmide quadrangular
  \begin{scope}[xshift=13cm, yshift=1.5cm]
    \draw (0,0) rectangle (1.5,1.5);
    \draw (0,1.5) -- (0.75,3) -- (1.5,1.5);
    \draw (0,0) -- (0.75,-1.5) -- (1.5,0);
    \draw (0,0) -- (-1.5,0.75) -- (0,1.5);
    \draw (1.5,0) -- (3,0.75) -- (1.5,1.5);
    \node at (0.75,-2.3) {(c)};
  \end{scope}
  % (d) cilindro
  \begin{scope}[xshift=19cm]
    \draw (0,0.5) rectangle (3.2,2.5);
    \draw (1.6,3.0) circle (0.5);
    \draw (1.6,0) circle (0.5);
    \node at (1.6,-0.8) {(d)};
  \end{scope}
\end{tikzpicture}
\end{center}

\begin{exercicios}
  \item Identifica os sólidos das planificações (a), (b), (c) e (d).
  \item Na planificação (a), que face fica oposta à face do topo?
  \item Na planificação (d), qual é a relação entre o comprimento do retângulo e o raio dos círculos?
  \item Desenha uma planificação de um paralelepípedo com 4\,cm $\times$ 2\,cm $\times$ 1\,cm.
  \item Desenha uma planificação de um cone.
  \item Quantas planificações diferentes tem um cubo? (Tenta encontrar pelo menos 5.)
\end{exercicios}

\secao{Vistas e Secções}{Responde a cada questão.}
\begin{exercicios}
  \item Que figura se obtém ao cortar um cilindro por um plano paralelo às bases?
  \item Que figura se obtém ao cortar um cone por um plano paralelo à base?
  \item Que figura se obtém ao cortar um cubo por um plano que passa por duas arestas opostas?
  \item Que figura se obtém ao cortar uma esfera por um plano?
  \item Desenha a vista de frente, de cima e de lado de uma pirâmide quadrangular.
  \item Um sólido tem vista de frente e de lado triangulares, e vista de cima circular. Que sólido é?
\end{exercicios}

\secao{Desafios}{Resolve os problemas.}
\begin{exercicios}
  \item Os cinco poliedros regulares (sólidos platónicos) são o tetraedro, o cubo, o octaedro, o dodecaedro e o icosaedro. Indica o tipo de face de cada um.
  \item Um icosaedro tem 20 faces triangulares. Usando a relação de Euler, determina o número de arestas e de vértices.
  \item Um cubo de aresta 3\,cm, pintado de vermelho, é cortado em cubinhos de 1\,cm. Quantos cubinhos têm 3 faces pintadas? E 2? E 1? E nenhuma?
  \item Uma pirâmide e um prisma têm a mesma base com $n$ lados. Juntos têm 25 faces. Determina $n$.
  \item Explica por que razão não existe nenhum poliedro com exatamente 7 arestas. (Sugestão: cada face tem pelo menos 3 arestas e cada vértice pelo menos 3 arestas, logo $3F \le 2A$ e $3V \le 2A$.)
\end{exercicios}

\end{document}
